\section{Aplicaciones e impacto del Transformer}

\subsection{Attention Is All You Need}

En 2017, Vaswani et al. publicaron el artículo histórico \textit{Attention Is All You Need}, en el que se propuso la arquitectura Transformer. Esta arquitectura \textbf{prescinde por completo del mecanismo recurrente} y construye el modelo de secuencias basándose únicamente en Self-Attention y capas de retroalimentación directa.

\begin{figure}[H]
\centering
\includegraphics[width=0.6\textwidth]{slides-images/slide-65.jpg}
\caption{«Attention Is All You Need»: el inicio de la era del Transformer\protect\footnotemark}
\end{figure}
\footnotetext{Slides página 65.}

\begin{knowledgebox}{La T de GPT significa Transformer}
Si has usado ChatGPT, entonces la «T» de su nombre representa Transformer. GPT significa Generative Pre-trained Transformer, es decir, un modelo generativo preentrenado basado en la arquitectura Transformer. El Transformer es hoy la arquitectura subyacente de casi todos los grandes modelos de lenguaje.
\end{knowledgebox}

\subsection{Transformer contra RNN: comparación de ventajas}

\begin{table}[H]
\centering
\caption{Comparación entre RNN y Transformer}
\begin{tabular}{lcc}
\toprule
\textbf{Característica} & \textbf{RNN/LSTM} & \textbf{Transformer} \\
\midrule
Modo de procesamiento & Paso a paso en el tiempo (secuencial) & Toda la secuencia en paralelo \\
Dependencias de largo plazo & Difícil (desvanecimiento del gradiente) & Modelado directo (Self-Attention) \\
Eficiencia de entrenamiento & Baja (no se puede paralelizar) & Alta (aprovecha bien la GPU) \\
Información de posición & Implícita (por el orden de los pasos de tiempo) & Explícita (Positional Encoding) \\
Escalabilidad & Limitada & Fuerte (admite escalas muy grandes) \\
\bottomrule
\end{tabular}
\end{table}

\subsection{El amplio alcance de Self-Attention}

El impacto de Self-Attention y del Transformer va mucho más allá del procesamiento de lenguaje natural:

\begin{figure}[H]
\centering
\includegraphics[width=0.85\textwidth]{slides-images/slide-80.jpg}
\caption{Aplicaciones de Self-Attention en distintos campos\protect\footnotemark}
\end{figure}
\footnotetext{Slides página 80.}

\begin{itemize}
    \item \textbf{Procesamiento de lenguaje natural}: grandes modelos de lenguaje como BERT (Devlin et al., NAACL 2019) y GPT (Brown et al., NeurIPS 2020)
    \item \textbf{Secuencias biológicas}: AlphaFold (Jumper et al., Nature 2021) usa Attention para predecir la estructura tridimensional de proteínas; ESMFold (Lin et al., Science 2023)
    \item \textbf{Visión por computadora}: Vision Transformer / ViT (Dosovitskiy et al., ICLR 2020) divide la imagen en una secuencia de parches (Patch) y la procesa con un Transformer, con un desempeño comparable o incluso superior al de las redes neuronales convolucionales
    \item \textbf{Multimodalidad}: los mecanismos de Attention entre modalidades permiten tareas como la generación de imágenes a partir de texto (como DALL-E) y la comprensión de video
\end{itemize}

\begin{importantbox}{La generalidad del Transformer}
La clave del éxito del Transformer está en su diseño arquitectónico \textbf{independiente del dominio} (Domain-agnostic). Sin importar si la entrada es texto, imágenes, secuencias de proteínas o audio, basta con poder representarla como una serie de vectores (Token) para poder procesarla con un Transformer. Esta generalidad ha hecho del Transformer la arquitectura estándar de la era de los «modelos fundacionales» (Foundation Model).
\end{importantbox}

\subsection{Resumen de la sección}

El Transformer transformó por completo el panorama del modelado de secuencias. A través del mecanismo de Self-Attention eliminó la dependencia de estructuras recurrentes, logrando cómputo paralelo eficiente y modelado directo de dependencias de largo alcance. Desde el procesamiento de lenguaje natural hasta la predicción de estructuras de proteínas, y desde la visión por computadora hasta la generación multimodal, el impacto del Transformer está presente en todas partes.

\newpage
